% !TeX root = ../main.tex

\chapter{复现一致性与全期表现}
\section{原复现区间的核对}
基础复现与延展在147个共同收益月份完全一致。四路径初始20个季度的256条股票—季度记录均与基础版逐股一致；P3、P4名单相同，但由于仓位管理，初始期收益不必与未控组合相同。

\begin{table}[htbp]
\centering
\caption{共同复现区间的绩效（2003-04至2015-06）}
\renewcommand{\arraystretch}{1.18}
\small
\setlength{\tabcolsep}{4.5pt}
\begin{tabularx}{\textwidth}{@{}p{2.55cm}*{4}{C}@{}}
\toprule
方案 & 月数 & 年化收益/\% & 波动/\% & 最大回撤/\% \\
\midrule
基础复现 B & 147 & 38.86 & 35.97 & -70.11 \\
基础延展 E & 147 & 38.86 & 35.97 & -70.11 \\
路径一 P1 & 147 & 41.73 & 35.85 & -68.87 \\
路径二 P2 & 147 & 43.33 & 36.14 & -64.15 \\
路径三 P3 & 147 & 30.61 & 23.48 & -32.74 \\
路径四 P4 & 147 & 30.92 & 24.39 & -34.26 \\
上证指数 & 147 & 9.09 & 28.93 & -70.97 \\
\bottomrule
\end{tabularx}
\end{table}

基础规则在共同区间的年化收益为38.86\%，最大回撤为$-$70.11\%。P2在该区间年化43.33\%，高于P1的41.73\%；后文可见，这一排序在完整延展区间逆转。因此，不能仅凭原复现终点选择所谓最优排名方式。本研究为月频近似，不把原研报的不同起止日期和日频结果当作逐日精确核对标准。

\begin{figure}[htbp]
\centering
\includegraphics[width=\linewidth]{fig02}
\caption{全部方案的共同区间净值}
\par\smallskip
\begin{minipage}{0.97\linewidth}\small 注：基础复现与延展共同月份完全重合，仅展示一条基础曲线；使用对数坐标。资料来源：CSMAR，作者计算。\end{minipage}
\end{figure}

\section{延展后的收益与风险}
\begin{table}[htbp]
\centering
\caption{完整延展区间的收益与风险（2003-04至2026-09）}
\renewcommand{\arraystretch}{1.18}
\small
\setlength{\tabcolsep}{4.5pt}
\begin{tabularx}{\textwidth}{@{}p{2.35cm}*{5}{C}@{}}
\toprule
方案 & 年化收益/\% & 波动/\% & 最大回撤/\% & 零利率夏普 & Calmar \\
\midrule
基础延展 E & 26.17 & 32.41 & -70.11 & 0.881 & 0.373 \\
路径一 P1 & 29.64 & 29.98 & -68.87 & 1.021 & 0.430 \\
路径二 P2 & 27.38 & 30.93 & -64.15 & 0.941 & 0.427 \\
路径三 P3 & 23.19 & 21.45 & -32.74 & 1.084 & 0.708 \\
路径四 P4 & 20.85 & 22.21 & -34.86 & 0.967 & 0.598 \\
上证指数 & 4.13 & 24.13 & -70.97 & 0.289 & 0.058 \\
\bottomrule
\end{tabularx}
\end{table}

282个月中，E年化26.17\%、波动32.41\%、最大回撤$-$70.11\%。P1年化29.64\%，较E增加3.48个百分点；波动降为29.98\%，但最大回撤仍达$-$68.87\%，仅收窄1.24个百分点。增加持股和提高夏普，并没有消除极端峰谷损失。

P2年化27.38\%、波动30.93\%、最大回撤$-$64.15\%。相较P1，它以2.26个百分点年化收益的下降，换取4.72个百分点的回撤收窄；波动反而上升0.95个百分点。故排名恢复没有同时改善收益、波动和回撤，也不能被概括为全面优于候选池内排名。

\begin{figure}[htbp]
\centering
\includegraphics[width=\linewidth]{fig03}
\caption{基础延展与两条选股优化路径的全期净值及回撤}
\par\smallskip
\begin{minipage}{0.97\linewidth}\small 注：上图为对数净值，下图为从初始净值1计算的月末回撤；同一区间282个月。资料来源：CSMAR，作者计算。\end{minipage}
\end{figure}

\section{仓位管理的成对影响}
P3将P1波动由29.98\%降至21.45\%，最大回撤由$-$68.87\%收窄至$-$32.74\%；代价是年化从29.64\%降至23.19\%。P4将P2波动由30.93\%降至22.21\%，最大回撤由$-$64.15\%收窄至$-$34.86\%，年化从27.38\%降至20.85\%。两组均体现风险暴露下降与上涨参与下降并存。

P3零利率夏普1.084，高于P1的1.021；P4为0.967，高于P2的0.941。风控改善了所设夏普，但没有提高全期绝对收益。P3在本次完整区间同时高于P4的收益、低于其波动和回撤；这只是指定规则及样本中的排序，并非对未来或所有风险偏好的支配结论。

\begin{figure}[htbp]
\centering
\includegraphics[width=\linewidth]{fig04}
\caption{两条风控路径的全期净值及回撤}
\par\smallskip
\begin{minipage}{0.97\linewidth}\small 注：基础延展不额外叠加风控，作为相同信息口径的原规则对照。资料来源：CSMAR，作者计算。\end{minipage}
\end{figure}

\begin{figure}[htbp]
\centering
\includegraphics[width=\linewidth]{fig05}
\caption{各正式路径的年化收益与最大回撤}
\par\smallskip
\begin{minipage}{0.97\linewidth}\small 注：斜线柱为动态仓位路径；最大回撤越接近零表示峰谷损失越小。资料来源：CSMAR，作者计算。\end{minipage}
\end{figure}

\section{如何理解收益提高与风险降低}
全期年化收益回答的是持有整个样本的复利速度，年化波动描述月收益离散程度，最大回撤描述历史峰值到后续谷值的最大损失。三者不是同一风险收益评价。P1相对E同时提高年化、降低波动，但最大回撤只改善1.24个百分点，说明常态收益分布有所改善仍可与极端损失集中出现并存。不能把波动下降或持股增加直接表述为历史大跌已被化解。

P1与P2的比较同样需要多指标阅读：P2全期收益较低、最大回撤较小，月收益波动却略高。因此，“风险更低”必须明确指哪项风险。若研究目标是提高绝对复利，当前全样本结果支持P1居前；若关注原始排名解释性，P2提供更接近基础版的母体定义；若主要控制本样本峰谷损失，则需要连同P3、P4的收益代价评价，不能把三种目标合并为单一“最好”。

累计财富差异会被长期复利放大。P1与P2的年化只相差约2.26个百分点，经过23.5年，终值比例仍可能有明显差距；这种放大不意味着每年都领先，也不意味着收益差具有统计显著性。后续年度和分段表负责检验领先是否集中于少数时期，而非根据全期终值给出稳定性结论。

\section{研报原始时期内的优化效果检验}
\subsection{固定截止范围与比较窗口}
此前共同147个月的结果以2015年6月结束，而原研报公布的结果至2015年7月31日。为检验优化是否在原始时期内同样改善，本节从CSMAR逐股月收益独立重算，统一保留2015年5月31日前的选股信号；按本项目季度末规则，最后信号为3月31日，共49次。末组持仓由4—6月延持至7月，不加入6月新信号，不在没有新交易的7月另扣一次季度费用。因此，本节没有用原资料截止后才形成的新季度信号填充最后一个月。

评价起止分为三种：147个月为2003年4月至2015年6月，与既有基础版对应；148个月保留相同起点并延长末组至7月；151个月在148个月前加入2003年1—3月建仓前现金。现金仅用于起始分母敏感性，不作为已确认的原研报建仓规则，也不给风控增加虚构的预热收益。上证指数在1—3月使用实际CSMAR收益，不跟随策略置为现金。

风控严格沿用“选股信号变化时重置参考净值”的原机制。延持7月没有新信号，故仍使用末组参考值；历史波动窗口继续按此前收益更新。成熟期从2008年4月至2015年7月，共88个月，控制路径继承先前状态，不在2008年重新启动。月末回撤在每个评价窗口包含初始净值1。

\begin{table}[htbp]
\centering
\caption{原始时期不同起止口径的年化收益}
\renewcommand{\arraystretch}{1.18}
\small
\setlength{\tabcolsep}{4.5pt}
\begin{tabularx}{\textwidth}{@{}p{2.55cm}*{4}{C}@{}}
\toprule
方案 & 147月/\% & 148月/\% & 151月/\% & 151月相对B/百分点 \\
\midrule
基础复现B & 38.86 & 37.31 & 36.44 & +0.00 \\
路径一P1 & 41.73 & 39.70 & 38.77 & +2.33 \\
路径二P2 & 43.33 & 41.95 & 40.97 & +4.52 \\
路径三P3 & 30.61 & 29.08 & 28.43 & -8.01 \\
路径四P4 & 30.92 & 30.14 & 29.46 & -6.98 \\
上证指数 & 9.09 & 7.67 & 8.42 & -28.02 \\
\bottomrule
\end{tabularx}
\par\smallskip
\begin{minipage}{\textwidth}\small 注：147月、148月、151月并非三次独立投资试验，而是同一原始时期的起止敏感性；相对B为同一151个月内的CAGR差。申万公布的36.74\%仍具有不同公告日调仓及日频定义，本表不把它当作已经同口径的基准。\end{minipage}
\end{table}

\subsection{收益改善及其风险代价}
三种窗口均出现P1、P2年化高于B，P3、P4年化低于B的方向。在151个月口径下，B年化36.44\%，P1为38.77\%、P2为40.97\%，分别提高2.33和4.52个百分点。将终点从6月推至7月后，这一选股改善方向仍保留，而不是仅在某一个截止月份成立。P2在原始时期高于P1，亦与后文完整282个月中P1领先的排序不同，进一步说明规则的相对表现依赖市场阶段。

P1、P2的最大回撤分别为$-$68.87\%、$-$64.15\%，相对B的$-$70.11\%收窄1.24、5.96个百分点。但年化波动与B接近或略高，所以“收益和最大回撤改善”不等于全部风险指标同步改善。尤其P1仍有接近69\%的历史峰谷损失，不能因年化增加就忽略损失承受要求。

P3、P4年化降至28.43\%、29.46\%，较B下降8.01、6.98个百分点；回撤则收窄至$-$32.74\%、$-$34.26\%，零利率夏普分别为1.187、1.195，高于B的1.055。这支持风控在原始时期降低风险暴露的描述，却不支持其提高全原始区间绝对收益。此处P4的$-$34.26\%与完整延展的$-$34.86\%来自不同终点，应分别阅读。

\begin{table}[htbp]
\centering
\caption{原始时期风险及成熟期的补充检验}
\renewcommand{\arraystretch}{1.18}
\small
\setlength{\tabcolsep}{4.5pt}
\begin{tabularx}{\textwidth}{@{}p{2.35cm}*{5}{C}@{}}
\toprule
方案 & 151月回撤/\% & 151月波动/\% & 151月夏普 & 成熟期年化/\% & 成熟期回撤/\% \\
\midrule
基础复现B & -70.11 & 35.74 & 1.055 & 18.45 & -62.72 \\
路径一P1 & -68.87 & 35.74 & 1.102 & 21.94 & -60.64 \\
路径二P2 & -64.15 & 35.86 & 1.143 & 25.26 & -54.33 \\
路径三P3 & -32.74 & 23.53 & 1.187 & 19.10 & -26.15 \\
路径四P4 & -34.26 & 24.18 & 1.195 & 20.75 & -34.26 \\
上证指数 & -70.97 & 28.99 & 0.425 & 1.08 & -53.19 \\
\bottomrule
\end{tabularx}
\par\smallskip
\begin{minipage}{\textwidth}\small 注：成熟期88个月为2008-04至2015-07；分段回撤按段初净值1重算，不能替代跨阶段回撤。夏普沿用零无风险利率口径，收益、回撤与波动为百分数。\end{minipage}
\end{table}

\subsection{排除共同初始阶段后如何判断}
成熟期88个月中，B年化18.45\%，P1、P2为21.94\%、25.26\%，仍分别提高3.49、6.81个百分点。因此，两条选股改进的优势并非只由完全相同的初始五年名单产生。P3、P4在该段年化19.10\%、20.75\%，也分别高于B约0.65、2.30个百分点，但仍低于各自底层P1、P2。

这说明“风控在整个原始时期降低收益”与“成熟期风控路径收益高于基础版”可以同时成立：前者与原始B比较并包括初期高涨，后者只评价成熟期，且底层选股也发生变化。不能据成熟期P3高于B，就把收益增加单独归于风控；识别仓位作用仍应比较P1与P3、P2与P4。初始高涨参与减少与阶段排序共同提供解释，但没有构成各机制可加总的因果分解。

本实验24行指标覆盖6个方案和4个窗口，R逐股重算后与原147个月逐月一致；Python再次独立连接CSMAR收益、重建148个月控制仓位并复算指标。其证据支持指定月频口径下的改善方向和计算一致性，尚未解决真实公告时点、原历史覆盖及成交约束，也不是事先冻结的样本外验证。完整实验程序和逐月结果独立保存在“研报原始区间优化验证\_20261107”。

\begin{figure}[htbp]
\centering
\includegraphics[width=\linewidth]{fig06}
\caption{原始时期各路径的收益与最大回撤}
\par\smallskip
\begin{minipage}{0.97\linewidth}\small 注：151个月；建仓前现金及末组延持仅为明示敏感性，不是申万精确日频重建。资料来源：CSMAR，作者计算。\end{minipage}
\end{figure}

\begin{figure}[htbp]
\centering
\includegraphics[width=\linewidth]{fig07}
\caption{原始时期的月度净值与回撤过程}
\par\smallskip
\begin{minipage}{0.97\linewidth}\small 注：实际投资期148个月；统一2015年5月前信号，不使用6月新信号。资料来源：CSMAR，作者计算。\end{minipage}
\end{figure}
